% Complete Section 8.6, original PDF 113-114 / printed 103-104.
% block: maxwell-equations-heading page:113 type:section
\subsection{{Maxwell's equations}}\label{sec:maxwell-equations}

% block: maxwell-traditional-system page:113 type:ordinary-paragraph
{In this section we show how to formulate Maxwell's equations as, essentially, a Hodge wave equation. The equations are, in traditional notation,}
\begin{equation}\label{eq:maxwell-ampere}D_t-\operatorname{curl}H=-j,\end{equation}
\begin{equation}\label{eq:maxwell-faraday}B_t+\operatorname{curl}E=0,\end{equation}
\begin{equation}\label{eq:maxwell-magnetic-divergence}\operatorname{div}B=0,\end{equation}
\begin{equation}\label{eq:maxwell-electric-divergence}\operatorname{div}D=q.\end{equation}

% block: maxwell-constitutive-form-interpretation page:113 type:ordinary-paragraph
{We view the electric field $E$ and the magnetizing field $H$ as (vector proxies for) 1-forms and the magnetic field $B$ and electric displacement $D$ as 2-forms. Thus, both the curl operators in the equations are the exterior derivative $d^1$ and both the divergence operators are $d^2$. The current density $j$ is thus a 2-form, and the charge density $q$ a 3-form. These equations are supplemented by the constitutive equations}
\[
D=\epsilon E,\qquad B=\mu H,
\]
{where the permittivity $\epsilon$ and permeability $\mu$ are positive scalars or symmetric positive definite operators which may vary in space (but not time). Note that equation $D=\epsilon E$ has a clear meaning if we view both $D$ and $E$ as vector fields but less so if we think of it as relating differential forms of different orders. Written correctly in the language of differential forms, the equation is $D=\star_\epsilon E$, where the Hodge star operator $\star_\epsilon$ is formed with respect to an $\epsilon$-weighted inner product rather than the usual unweighted $L^2\Lambda^k$ inner product.}

% block: maxwell-data-compatibility page:113 type:ordinary-paragraph
{We note that \eqref{eq:maxwell-ampere} and \eqref{eq:maxwell-electric-divergence} imply a necessary compatibility condition on the data, namely, that $q_t=-\operatorname{div}j$. In the notation of exterior calculus, this says that $q_t=-dj$ or, weakly,}
\begin{equation}\label{eq:maxwell-data-compatibility}
\int q_t\wedge\tau-\int j\wedge d\tau=0,\qquad\tau\in\mathring H\Lambda^0.
\end{equation}

% block: maxwell-initial-boundary-data page:113 type:ordinary-paragraph
{The initial data we impose on $E$, $D$, $B$, and $H$ must be compatible as well. Obviously we must have}
\[
D(0)=\epsilon E(0),\qquad B(0)=\mu H(0),
\]
{and, from \eqref{eq:maxwell-magnetic-divergence} and \eqref{eq:maxwell-electric-divergence}, we require}
\[
\operatorname{div}B(0)=0,\qquad\operatorname{div}D(0)=q(0),
\]
{as well. Finally the Maxwell problem is completed with boundary conditions. For simplicity, we take homogeneous electric (perfect conductor) conditions $E\times n=0$, $B\cdot n=0$, i.e., $\operatorname{tr}E=0$ and $\operatorname{tr}B=0$.}

% block: maxwell-slack-variable-and-weights pages:113,114 type:ordinary-paragraph
{Our immediate goal is to write Maxwell's equations in the form of the 1-form Hodge wave equation. The unknown time-dependent 1-form, denoted by $v$ in \eqref{eq:hodge-wave-weak-formulation}, will be the electric field $E$. The unknown 2-form, denoted by $\beta$ above, will be the negative of the magnetic field $B$. The remaining dependent variable, the 0-form $\sigma$, will in fact turn out to vanish identically. It is a slack variable, whose vanishing reflects the compatibility of the data $q$ and $j$. (Our weak formulation will be well-posed even when the data $q$ and $j$ are incompatible, but, in that case, $\sigma$ will not vanish, and $E$ and $B$ will not satisfy Maxwell's equations.) As explained above, the variable coefficients $\epsilon$ and $\mu$ are handled by adjusting the inner product we use on the spaces of $L^2$ differential forms. For differential 1-forms we weight the inner product by $\epsilon$,}
\[
\langle E,F\rangle_\epsilon:=\int\star_\epsilon E\wedge F,
\]
{which is simply $\int\epsilon E\cdot F\,dx$ when written in terms of the vector proxies. Similarly for differential 2-forms we weight the inner product by $\mu^{-1}$. For the 0-forms we may weight with an arbitrary positive function $\alpha$.}

% block: maxwell-weak-system page:114 type:ordinary-paragraph
{We are now ready to write the Maxwell system in weak form. With $\sigma\equiv0$, we get from \eqref{eq:maxwell-ampere}--\eqref{eq:maxwell-electric-divergence} that $(\sigma,E,B):[0,T]\to\mathring H\Lambda^0\times\mathring H\Lambda^1\times\mathring H\Lambda^2$ satisfy}
\begin{equation}\label{eq:maxwell-weak-slack}
\langle\sigma_t,\tau\rangle_\alpha-\langle E,d\tau\rangle_\epsilon=\int q\wedge\tau,\qquad\tau\in\mathring H\Lambda^0,\ t\in[0,T],
\end{equation}
\begin{equation}\label{eq:maxwell-weak-electric}
\langle E_t,F\rangle_\epsilon+\langle d\sigma,F\rangle_\epsilon-\langle B,dF\rangle_{\mu^{-1}}=-\int j\wedge F,\quad F\in\mathring H\Lambda^1,\ t\in[0,T],
\end{equation}
\begin{equation}\label{eq:maxwell-weak-magnetic}
\langle B_t,C\rangle_{\mu^{-1}}+\langle dE,C\rangle_{\mu^{-1}}=0,\qquad C\in\mathring H\Lambda^2,\ t\in[0,T],
\end{equation}
{with the initial conditions $\sigma(0)=0$, $E(0)=E_0$, $B(0)=B_0$. After changing variables from $B$ to $-B$, this problem is exactly of the form of the Hodge wave equation \eqref{eq:hodge-wave-weak-formulation}, except that the inner products are weighted, and the right-hand side of the first equation need not vanish. The well-posedness of the problem can then be established in the same way as for \eqref{eq:hodge-wave-weak-formulation}.}

% block: maxwell-slack-variable-vanishing page:114 type:ordinary-paragraph
{Let us show that, under the assumption of compatible data, the slack variable $\sigma$ vanishes and that, with $D=\star_\epsilon E$ and $H=\star_{\mu=1}B$, the quadruple $(E,B,H,D)$ satisfies Maxwell's equations. To see that $\sigma$ vanishes, we first note that it vanishes at time $t=0$ by assumption. Next, we evaluate \eqref{eq:maxwell-weak-slack} at time $t=0$ and use the fact that $\star_\epsilon E(0)=D(0)$ and $dD(0)=q(0)$ to conclude that $\sigma_t(0)=0$. Next, we differentiate \eqref{eq:maxwell-weak-slack} in time and add it to \eqref{eq:maxwell-weak-electric}, taking $F=d\tau$ in the latter. Using the compatibility condition \eqref{eq:maxwell-data-compatibility}, we obtain}
\[
\langle\sigma_{tt},\tau\rangle_\alpha+\langle d\sigma,d\tau\rangle_\epsilon=0,\qquad\tau\in\mathring H\Lambda^0,\ t\in[0,T].
\]
{Taking $\tau=\sigma_t$, we obtain that}
\[
\frac d{dt}(\|\sigma_t\|_\alpha^2+\|d\sigma\|_\epsilon^2)=0.
\]
{Thus the quantity in parentheses is constant in time, and we have already seen that it vanishes when $t=0$. Therefore, $\sigma_t$ vanishes, i.e., $\sigma$ is constant in time, and so it too must vanish.}

% block: maxwell-equations-recovered page:114 type:ordinary-paragraph
{Having established that $\sigma$ vanishes, it is easy to see that \eqref{eq:maxwell-weak-slack}--\eqref{eq:maxwell-weak-magnetic} imply the Maxwell equations. Equation~\eqref{eq:maxwell-weak-slack} is a weak restatement of \eqref{eq:maxwell-electric-divergence}, \eqref{eq:maxwell-weak-electric} of \eqref{eq:maxwell-ampere}, and \eqref{eq:maxwell-weak-magnetic} of \eqref{eq:maxwell-faraday}. Finally \eqref{eq:maxwell-magnetic-divergence} follows from \eqref{eq:maxwell-faraday} since $\operatorname{div}B(0)=0$.}
